\documentclass[a4paper,12pt]{article}

\usepackage[utf8]{inputenc}
\usepackage[T2A]{fontenc}
\usepackage[russian]{babel}
\usepackage{amsmath}
\usepackage{amssymb}
\usepackage{graphicx}
\usepackage{geometry}
\usepackage{hyperref}
\usepackage{float}

\geometry{margin=2.5cm}

\hypersetup{
    pdfborderstyle={/S/U/W 1},	
    colorlinks=true,
    linkcolor=blue,
    filecolor=magenta,      
    urlcolor=blue,
    pdfpagemode=FullScreen,
    }

\title{Лабораторная работа \\ 
Дополненная реальность: оценка позы ArUco-маркеров и построение 3D-сцены}
\author{}
\date{}

\begin{document}

\maketitle


\section{Введение}

Дополненная реальность (Augmented Reality, AR) представляет собой класс систем, в которых виртуальные объекты интегрируются в изображение реальной сцены таким образом, чтобы сохранялась геометрическая согласованность между реальными и синтетическими объектами. Ключевой задачей любой AR-системы является оценка положения и ориентации камеры относительно наблюдаемой сцены, то есть задача оценки позы (pose estimation).

В общем случае определение позы камеры может выполняться по естественным особенностям изображения (feature-based методы), по инерциальным данным (различным датчикам) или по заранее известным объектам в сцене. Однако методы, основанные на естественных признаках, обладают рядом ограничений: чувствительность к текстуре, освещению, повторяющимся структурам, а также сложность обеспечения высокой точности в реальном времени.

Для повышения устойчивости и точности в системах дополненной реальности широко используются искусственные визуальные маркеры. Маркер представляет собой специально сконструированный объект с известной геометрией и кодированной структурой, позволяющей:

\begin{itemize}
    \item однозначно идентифицировать объект в кадре;
    \item получить устойчивые 2D--3D соответствия;
    \item вычислить положение и ориентацию маркера относительно камеры;
    \item задать локальную систему координат сцены.
\end{itemize}

В отличие от произвольных текстур, маркеры проектируются таким образом, чтобы их детекция была устойчива к изменениям масштаба, поворота и перспективных искажений. Наличие жёстко заданной геометрии (как правило, квадратной формы) существенно упрощает решение задачи оценки позы, поскольку координаты характерных точек маркера в собственной системе координат известны заранее.

Использование маркеров в AR-системах позволяет перейти от абстрактной задачи реконструкции сцены к более строго формализованной задаче определения жёсткого преобразования между системой координат камеры и системой координат маркера. После определения этого преобразования становится возможным корректное размещение виртуальных объектов в реальном пространстве с сохранением перспективы, масштаба и ориентации.

Таким образом, маркер выступает в роли геометрического «якоря» сцены: он фиксирует положение виртуальной системы координат в физическом мире и обеспечивает воспроизводимость пространственной конфигурации. В рамках данной лабораторной работы маркеры используются как опорные элементы для построения общей системы координат стола и размещения виртуальной трёхмерной сцены.




\section{Постановка задачи}

На физическом столе закреплены 2--4 ArUco-маркера. Камера наблюдает сцену в реальном времени. Параметры калибровки камеры считаются известными.

Необходимо разработать систему дополненной реальности, которая:

\begin{itemize}
    \item детектирует маркеры в реальном времени;
    \item оценивает их позу;
    \item формирует систему координат стола;
    \item строит виртуальную пирамиду из кубиков;
    \item корректно размещает пирамиду на плоскости стола.
\end{itemize}

Архитектура решения полностью определяется студентом.

\section{Системы координат и жёсткие преобразования}

\setcounter{equation}{0}

\subsection{Используемые системы координат}

Рассмотрим следующие системы координат:

\begin{itemize}
    \item мировая система $\mathcal{W}$;
    \item система камеры $\mathcal{C}$;
    \item система маркера $\mathcal{M}$.
\end{itemize}

Каждая система координат является декартовой трёхмерной системой с ортонормированным базисом.

Пусть $X_{\mathcal{W}} \in \mathbb{R}^3$ --- координаты точки в мировой системе координат,  
$X_{\mathcal{C}} \in \mathbb{R}^3$ --- координаты той же точки в системе камеры.

Связь между ними задаётся жёстким преобразованием.

\subsection{Жёсткое преобразование в трёхмерном пространстве}

Жёстким (евклидовым) преобразованием называется отображение, сохраняющее расстояния и углы. В $\mathbb{R}^3$ оно состоит из поворота и переноса.

Пусть:
\[
R \in SO(3), \qquad t \in \mathbb{R}^3.
\]

Тогда преобразование из системы $\mathcal{W}$ в систему $\mathcal{C}$ записывается как

\begin{equation}
X_{\mathcal{C}} = R X_{\mathcal{W}} + t,
\end{equation}

где R --- матрица поворота, t --- вектор переноса.

\bigskip

Матрица поворота принадлежит специальной ортогональной группе:

\begin{equation}
R^T R = I,
\end{equation}

\begin{equation}
\det R = 1.
\end{equation}

Из ортогональности следует:

\begin{equation}
R^{-1} = R^T.
\end{equation}

Это свойство существенно упрощает обращение жёсткого преобразования.

\subsection{Однородное представление}

Для удобства композиции преобразований используется однородная форма записи.

Точка в однородных координатах:

\begin{equation}
\tilde{X} =
\begin{pmatrix}
X \\
1
\end{pmatrix}.
\end{equation}

Жёсткое преобразование записывается как матрица $4 \times 4$:

\begin{equation}
T =
\begin{pmatrix}
R & t \\
0 & 1
\end{pmatrix}.
\end{equation}

Тогда преобразование принимает вид:

\begin{equation}
\tilde{X}_{\mathcal{C}} = T \tilde{X}_{\mathcal{W}}.
\end{equation}

Поскольку $R^{-1} = R^T$, обратное преобразование вычисляется аналитически:

\begin{equation}
T^{-1} =
\begin{pmatrix}
R^T & -R^T t \\
0 & 1
\end{pmatrix}.
\end{equation}

Это выражение играет важную роль при переходе между системами координат маркера, камеры и стола.

\subsection{Композиция преобразований}

Пусть заданы два преобразования:

\begin{equation}
T_{1} =
\begin{pmatrix}
R_1 & t_1 \\
0 & 1
\end{pmatrix},
\qquad
T_{2} =
\begin{pmatrix}
R_2 & t_2 \\
0 & 1
\end{pmatrix}.
\end{equation}

Их композиция определяется матричным произведением:

\begin{equation}
T_{12} = T_2 T_1.
\end{equation}

В развернутом виде:

\begin{equation}
T_{12} =
\begin{pmatrix}
R_2 R_1 & R_2 t_1 + t_2 \\
0 & 1
\end{pmatrix}.
\end{equation}

Из этого выражения видно, что:

\begin{itemize}
    \item повороты перемножаются,
    \item перенос второго преобразования добавляется после поворота первого.
\end{itemize}

\subsection{Цепочки преобразований}

Если точка задана в системе маркера $\mathcal{M}$, а требуется получить её координаты в системе камеры $\mathcal{C}$ через мировую систему $\mathcal{W}$, используется композиция:

\begin{equation}
T_{\mathcal{C}\mathcal{M}} =
T_{\mathcal{C}\mathcal{W}} \,
T_{\mathcal{W}\mathcal{M}}.
\end{equation}

Тогда:

\begin{equation}
\tilde{X}_{\mathcal{C}} =
T_{\mathcal{C}\mathcal{M}} \,
\tilde{X}_{\mathcal{M}}.
\end{equation}

Именно такие цепочки преобразований координат лежат в основе размещения виртуальных объектов в сцене дополненной реальности.

\subsection{Геометрический смысл}

Жёсткое преобразование полностью определяет:

\begin{itemize}
    \item положение начала координат одной системы относительно другой;
    \item ориентацию базисных векторов;
    \item относительное расположение всех точек.
\end{itemize}

В задачах дополненной реальности оценка позы маркера фактически означает вычисление матрицы $T_{\mathcal{C}\mathcal{M}}$. После её получения любая точка виртуального объекта, заданная в системе маркера или стола, может быть корректно перенесена в систему камеры и затем спроецирована на изображение.









\section{ArUco-маркеры и процесс их детекции}

ArUco-маркер представляет собой квадратный бинарный шаблон с кодированной внутренней структурой и чётко выделенной границей. Геометрия маркера (квадрат) и дискретная структура кода позволяют получать устойчивые 2D--3D соответствия, необходимые для решения задачи оценки позы.

Каждый маркер принадлежит некоторому словарю. Словарь определяет:
\begin{itemize}
    \item размер внутренней бинарной сетки ($4\times4$, $5\times5$, $6\times6$ и т.д.);
    \item количество уникальных идентификаторов;
    \item минимальное кодовое расстояние между маркерами.
\end{itemize}

Чем больше размер сетки, тем выше потенциальная устойчивость к ошибкам и больше количество уникальных ID, однако усложняется детекция на больших расстояниях.

Процесс детекции маркера состоит из последовательности этапов:

\begin{enumerate}
    \item Предварительная обработка и бинаризация изображения;
    \item Поиск квадратных контуров;
    \item Проверка геометрических ограничений;
    \item Перспективное выравнивание (warp);
    \item Декодирование бинарного шаблона;
    \item Проверка словаря и коррекция ошибок.
\end{enumerate}

\subsection{Бинаризация изображения}

Цель первого этапа --- выделить потенциальные контуры маркеров.

Исходное изображение переводится в градации серого:

\begin{equation*}
\texttt{gray = cv2.cvtColor(frame, cv2.COLOR\_BGR2GRAY))}
\end{equation*}


Далее применяется адаптивная или глобальная пороговая обработка:
\begin{equation*}
\texttt{thresh = cv2.adaptiveThreshold(...)}
\end{equation*}

Адаптивная бинаризация предпочтительна, поскольку она компенсирует неравномерное освещение, повышает устойчивость к теням, облегчает выделение границ маркера.

Результатом является бинарное изображение, на котором потенциальные маркеры выделяются как контрастные квадраты.

\subsection{Поиск контуров и геометрическая фильтрация}
\bigskip

На бинарном изображении выполняется поиск контуров.

Для каждого контура проверяется замкнутость, площадь (фильтрация слишком маленьких или больших областей), приближение к четырёхугольнику.

После выполнения аппроксимации, например методом Дугласа–Пекера, как кандидаты принимаются те контуры, что:
\begin{itemize}
    \item содержат ровно 4 вершины;
    \item являются выпуклым;
    \item имеют разумные пропорции сторон.

    \item имеют корректную ориентация вершин (порядок обхода);
    \item не содержат самопересечений;
    \item имеют допустимое перспективное искажение.
\end{itemize}

Цель --- оставить только те четырёхугольники, которые могут быть перспективной проекцией квадрата.

\subsection{Перспективное выравнивание}

Для каждого кандидата вычисляется гомография, переводящая найденный четырёхугольник в нормализованный квадрат фиксированного размера.

Пусть заданы четыре точки изображения $x_i$ и соответствующие точки канонического квадрата $X_{\mathcal{M},i}$. Тогда существует матрица гомографии $H$, такая что:

\begin{equation}
x_i \sim H X_{\mathcal{M},i}.
\end{equation}

После вычисления гомографии выполняется перспективное преобразование (warp):

\begin{equation*}
\texttt{marker = cv2.warpPerspective(gray, H, (N, N))}
\end{equation*}

Результат --- изображение маркера в нормализованной системе координат без перспективных искажений.

\subsection{Бинарное декодирование}

Нормализованное изображение разбивается на сетку $n \times n$, соответствующую структуре словаря.

Для каждой ячейки вычисляется средняя интенсивность, и принимается решение:

\[
b_{ij} =
\begin{cases}
1, & \text{если интенсивность ниже порога}, \\
0, & \text{иначе}.
\end{cases}
\]

Полученная бинарная матрица интерпретируется как код маркера.

Важно учитывать возможные повороты маркера кратные $\frac{ \pi}{2}$ . Для каждого поворота формируется соответствующий код и сравнивается со словарём.

\subsection{Проверка словаря и коррекция ошибок}

Каждый словарь характеризуется минимальным расстоянием Хэмминга между кодами.

Расстояние Хэмминга определяется как:

\begin{equation}
d_H(a,b) = \sum_{i} |a_i - b_i|.
\end{equation}

Если расстояние до ближайшего допустимого кода меньше допустимого порога, маркер считается распознанным. Наличие такого кодового расстояния позволяет корректировать отдельные битовые ошибки и снижать вероятность ложных срабатываний.

\subsection{Формирование результата детекции}

После успешного декодирования формируются:

\begin{itemize}
    \item координаты четырёх углов маркера в изображении;
    \item идентификатор маркера (ID);
    \item порядок вершин.
\end{itemize}

В OpenCV итоговый вызов имеет вид:

\begin{equation*}
\texttt{corners, ids, rejected = cv2.aruco.detectMarkers(frame, dictionary)}
\end{equation*}

Полученные координаты углов служат входными данными для решения задачи PnP и последующей оценки позы.

\bigskip

При этом важно помнить, что данный процесс детекции чувствителен к сильным перспективным и нелинейным искажениям, малому размеру и частичной заслонённости маркера. При уменьшении видимого размера маркера уменьшается количество пикселей, соответствующих одной ячейке сетки, что приводит к росту ошибок бинаризации и декодирования.

Таким образом, ArUco-маркеры представляют собой специально сконструированный компромисс между информационной ёмкостью, устойчивостью к шуму и вычислительной эффективностью детекции.













\section{Оценка позы и задача Perspective-n-Point}

Оценка позы (pose estimation) --- ключевой этап построения системы дополненной реальности. Под оценкой позы понимается определение положения и ориентации объекта (или камеры) относительно другой системы координат.

В контексте данной лабораторной работы требуется определить жёсткое преобразование между системой координат маркера и системой координат камеры. Математически эта задача сводится к классической задаче \textbf{Perspective-n-Point (PnP)}.

\subsection{Постановка задачи PnP}

Пусть заданы:

\begin{itemize}
    \item набор трёхмерных точек $X_{\mathcal{M},i} \in \mathbb{R}^3$ в системе координат объекта;
    \item соответствующие им точки изображения $x_i \in \mathbb{R}^2$;
    \item матрица внутренних параметров камеры $K$.
\end{itemize}

Требуется определить матрицу поворота $R$ и вектор переноса $t$, такие что выполняется соотношение проекции:

\begin{equation}
x_i \sim K (R X_{\mathcal{M},i} + t)
\label{eq:pnp}
\end{equation}

Знак ``$\sim$'' означает равенство с точностью до масштаба в проективном пространстве.

Иначе говоря, необходимо найти такое жёсткое преобразование, при котором проекции трёхмерных точек максимально согласуются с наблюдаемыми точками изображения.

Для решения задачи необходимо минимум четыре некомпланарные точки.

В случае ArUco-маркера используются четыре угла квадрата, координаты которых известны в локальной системе координат маркера. Это делает маркер естественным объектом для постановки задачи PnP.

\subsection{Ошибка репроекции}

Поскольку измерения содержат шум, задача решается как задача оптимизации. Определяется ошибка репроекции:

\begin{equation}
e_i = x_i - \hat{x}_i(R,t),
\end{equation}

где $\hat{x}_i(R,t)$ --- проекция точки $X_i$ при текущих параметрах.

Суммарная квадратичная ошибка:

\begin{equation}
E(R,t) = \sum_{i=1}^{n} \| e_i \|^2.
\end{equation}

Задача PnP сводится к минимизации функционала (\theequation) по переменным $R$ и $t$.

\subsection{Параметризация поворота}

Матрица поворота имеет 9 элементов, но лишь 3 степени свободы. В численных алгоритмах поворот удобно задавать вектором Родригеса $r \in \mathbb{R}^3$.

Связь между вектором Родригеса и матрицей поворота определяется экспоненциальным отображением:

\begin{equation}
R = \exp([\omega]_\times),
\end{equation}

где $[\omega]_\times$ --- кососимметричная матрица, соответствующая вектору угловой скорости. В OpenCV данный переход выполняется следующим образом:

\begin{equation*}
\texttt{R, \_ = cv2.Rodrigues(rvec)}
\end{equation*}

Таким образом, оптимизация ведётся по параметрам $rvec$ и $tvec$.

\subsection{Алгоритмы решения PnP}

Существует несколько классов алгоритмов:

\begin{itemize}
    \item аналитические (P3P);
    \item линейные приближения (EPnP);
    \item итерационные методы (Levenberg--Marquardt).
\end{itemize}

Итерационные методы минимизируют ошибки репроекции напрямую и обеспечивают высокую точность при хорошем начальном приближении.

В OpenCV используется универсальный интерфейс:

\begin{equation*}
\texttt{retval, rvec, tvec = cv2.solvePnP(objPts, imgPts, K, dist)}
\end{equation*}

где:
\begin{itemize}
    \item \texttt{objPts} --- массив 3D-точек;
    \item \texttt{imgPts} --- соответствующие 2D-точки;
    \item \texttt{K} --- матрица внутренних параметров;
    \item \texttt{dist} --- коэффициенты дисторсии.
\end{itemize}

\subsection{Геометрическая интерпретация результата}

Результат решения PnP:

\begin{itemize}
    \item $rvec$ --- ориентация объекта относительно камеры;
    \item $tvec$ --- положение начала системы объекта в системе камеры.
\end{itemize}

Вектор переноса имеет прямой геометрический смысл:

\begin{equation}
tvec =
\begin{pmatrix}
t_x \\
t_y \\
t_z
\end{pmatrix},
\end{equation}

где $t_z$ соответствует расстоянию до объекта вдоль оптической оси камеры.

Матрица поворота определяет ориентацию локальных осей объекта в системе камеры.

Итоговое жёсткое преобразование:

\begin{equation}
T_{\mathcal{C}\mathcal{M}} =
\begin{pmatrix}
R & t \\
0 & 1
\end{pmatrix}
\label{eq:final_transform}
\end{equation}

\bigskip

Однако важно отметить, что задача PnP может иметь несколько решений (в случае малого числа точек) и является чувствительной к шуму, особенно возрастающему при малых углах наблюдения.

\bigskip

После получения $R$ и $t$ любая точка виртуального объекта, заданная в системе координат маркера, может быть преобразована в систему камеры, а затем спроецирована на изображение.
Таким образом, задача PnP является центральным вычислительным механизмом, обеспечивающим геометрическую согласованность виртуальной и реальной сцены.













\section{Проекция 3D-объектов на изображение}

После оценки позы объекта (маркера или стола) в системе координат камеры следующим этапом является отображение виртуальных трёхмерных объектов на изображение. Геометрически эта операция состоит из двух последовательных преобразований:

\begin{enumerate}
    \item переход из системы координат объекта в систему координат камеры;
    \item перспективная проекция точки из пространства $\mathbb{R}^3$ на плоскость изображения.
\end{enumerate}

\subsection{Переход из системы объекта в систему камеры}

Пусть точка виртуального объекта задана в системе координат $\mathcal{M}$:

\[
X_{\mathcal{M}} =
\begin{pmatrix}
X \\
Y \\
Z
\end{pmatrix}.
\]
Результатом решения задачи PnP является матрица поворота $R$ и вектор переноса $t$, определяющие преобразование из системы объекта в систему камеры:

\begin{equation}
X_{\mathcal{C}} = R X_{\mathcal{M}} + t.
\label{eq:xcxm}
\end{equation}

Здесь
\[
X_{\mathcal{C}} =
\begin{pmatrix}
X_c \\
Y_c \\
Z_c
\end{pmatrix}
\]
--- координаты точки в системе камеры.

\subsection{Перспективная проекция}

В модели pinhole-камеры точка проецируется на плоскость изображения по закону перспективы:

\begin{equation}
x' = \frac{X_c}{Z_c}, \qquad
y' = \frac{Y_c}{Z_c}.
\label{eq:startmap}
\end{equation}

Это нормализованные координаты в плоскости, расположенной на расстоянии $f=1$ от центра проекции.

С учётом реальных параметров камеры используется матрица внутренних параметров:

\begin{equation}
K =
\begin{pmatrix}
f_x & 0 & c_x \\
0 & f_y & c_y \\
0 & 0 & 1
\end{pmatrix}.
\end{equation}

Тогда координаты пикселя вычисляются как

\begin{equation}
\begin{pmatrix}
u \\
v \\
1
\end{pmatrix}
=
K
\begin{pmatrix}
x' \\
y' \\
1
\end{pmatrix}.
\end{equation}

Подставляя \eqref{eq:startmap}, получаем итоговую формулу проекции:

\begin{equation}
\begin{pmatrix}
u \\
v \\
1
\end{pmatrix}
=
K
\begin{pmatrix}
\frac{X_c}{Z_c} \\
\frac{Y_c}{Z_c} \\
1
\end{pmatrix}.
\label{eq:fullmap}
\end{equation}

В развернутом виде:

\begin{equation}
u = f_x \frac{X_c}{Z_c} + c_x,
\end{equation}

\begin{equation}
v = f_y \frac{Y_c}{Z_c} + c_y.
\end{equation}

Таким образом, размер объекта на изображении обратно пропорционален его глубине $Z_c$.

\subsection{Полная формула проекции}

Объединяя \eqref{eq:xcxm} и \eqref{eq:fullmap}, получаем компактную запись:

\begin{equation}
x \sim K (R X_{\mathcal{M}} + t).
\end{equation}

Именно это выражение лежит в основе функции \texttt{projectPoints}.

\begin{equation*}
\texttt{imgPts, \_ = cv2.projectPoints(points3D, rvec, tvec, K, dist)}
\end{equation*}

Функция выполняет:

\begin{itemize}
    \item преобразование координат через $R$ и $t$;
    \item деление на глубину $Z_c$;
    \item применение матрицы внутренних параметров;
    \item учёт дисторсии (если задана).
\end{itemize}

\subsection{Построение куба в локальной системе координат}

Пусть куб имеет ребро длины $L$ и его нижняя грань лежит в плоскости $Z=0$ системы координат объекта.

Вершины куба:

\begin{equation}
\begin{aligned}
&(0,0,0), (L,0,0), (L,L,0), (0,L,0), \\
&(0,0,L), (L,0,L), (L,L,L), (0,L,L)
\end{aligned}
\label{eq:cube}
\end{equation}

Каждая вершина рассматривается как точка $X_{\mathcal{M}}$ и подставляется в формулу \eqref{eq:pnp}.

\subsection{Алгоритм переноса куба на изображение}

Процесс отображения куба состоит из последовательных шагов:

\begin{enumerate}
    \item Задать массив из 8 вершин куба в системе координат объекта.
    \item Преобразовать каждую вершину в систему камеры.
    \item Выполнить перспективную проекцию.
    \item Соединить соответствующие пары точек рёбрами.
\end{enumerate}

После получения двумерных координат:

\begin{equation*}
\texttt{cv2.line(frame, pt1, pt2, thickness)}
\end{equation*}

соединяются рёбра:

\begin{itemize}
    \item четыре нижних ребра,
    \item четыре верхних,
    \item четыре вертикальных.
\end{itemize}

\subsection{Геометрический смысл}

Если преобразование $R,t$ оценено корректно, то:

\begin{itemize}
    \item куб будет жёстко «привязан» к маркеру;
    \item при движении камеры перспектива будет изменяться естественным образом;
    \item размер куба будет уменьшаться при увеличении расстояния;
    \item параллельные рёбра будут сходиться в точках схода.
\end{itemize}

Таким образом, перенос куба на изображение является прямым следствием последовательного применения жёсткого преобразования и перспективной проекции. Именно эта цепочка операций обеспечивает пространственную согласованность виртуального и реального мира в системе дополненной реальности.









\section{Техническое задание}

Система должна:

\begin{itemize}
    \item работать в реальном времени;
    \item использовать не менее двух маркеров;
    \item корректно размещать виртуальную пирамиду на столе;
    \item сохранять геометрическую согласованность при движении камеры.
\end{itemize}

\section{Рекомендуемые материалы}

Документация OpenCV (разделы aruco и calib3d):

\url{https://docs.opencv.org}

Материалы по задаче PnP и оценке позы в компьютерном зрении.

\end{document}